\section{Finite-window quantum spectra and magnetic crossover}
\label{sec:quantum-spectra}

The Weyl theorem fixes the smooth count but says nothing about local spectral
fluctuations.  We therefore performed a zero-input finite-window screen of the
centered \(a=1.02,n=1\) operator and its magnetic deformation.  The outcome
is a convergence-supported orthogonal-like to unitary-like response in the
available window.  It is not a proof of random-matrix universality, an
arithmetic selection of \(B\), or evidence for any individual Riemann zero.

\subsection{Discretization and frozen observables}

The primary discretization uses a Cartesian Dirichlet grid and Peierls links
\begin{equation}
 U_{ij}=\exp\!\left(-\ii\int_{q_i}^{q_j}A_B\cdot\dd q\right)
 \label{eq:peierls-link}
\end{equation}
on each nearest-neighbor edge.  Straight-link integration is exact for the
linear symmetric and Landau gauges used here.  The rectangle encloses the
preimage of the contour \(V=100E_{\mathrm{target}}\), with three nominal mesh
steps of padding; \(E_{\mathrm{target}}=450\), so the nominal wall energy is
\(45000\).  We computed 180 low eigenvalues by deterministic shift-invert
Hermitian Lanczos.  Every reported local statistic uses sorted modes
25--164 inclusive: 140 levels, 139 gaps, and 138 adjacent-gap ratios.

For consecutive gaps \(s_j=E_{j+1}-E_j\), the unfolding-free observable is
\begin{equation}
 \widetilde r_j=
 \frac{\min(s_j,s_{j+1})}{\max(s_j,s_{j+1})},
 \qquad
 \meanratio=\frac{1}{138}\sum_j\widetilde r_j.
 \label{eq:adjacent-ratio}
\end{equation}
The reference means used descriptively are \(0.38629\) for Poisson,
\(0.53590\) for GOE, and \(0.60266\) for GUE
\citep{AtasEtAl2013}.  We also unfolded with the
analytic clock
\begin{equation}
 \xi_j=\Ncl(E_j)
 =\frac{E_j}{2\pi}\log\frac{E_j}{2\pi}
  -\frac{E_j}{2\pi}+1.
 \label{eq:exact-clock-unfolding}
\end{equation}
No fit is involved in this transformation.  Because adjacent ratios from one
spectrum are dependent and the 140-level window was not sampled as an
i.i.d. ensemble, empirical CDF sup distances below are deterministic
discrepancy measures; we assign no \(p\)-values.

\subsection{The failed first grid comparison}

The first production run, R100, compared nominal spacings \(h=0.04\) and
\(0.03\).  It failed the frozen one-percent median level-change gate in every
cell: on the unified 25--164 window the changes ranged from \(1.20\%\) for
the radial control to \(1.67\%\) for the \(a=6\) cells, with the
\(a=1.02\) scalar and magnetic cells at \(1.23\%\) and \(1.21\%\).
Consequently, R100 alone supports no
random-matrix-class interpretation.

R101 was specified after that failure and added \(h=0.0225\).  Relative to
\(h=0.03\), all five cells then passed the original level gate, with median
changes between \(0.66\%\) and \(0.93\%\).  A post-hoc check compared each
fine-grid mean ratio with a two-grid \(h^2\) extrapolation.  The differences
were \(0.0120\) and \(0.0045\) for \(a=1.02\) at \(B=0\) and \(1\), but
\(0.0404\) and \(0.0347\) for \(a=6\).  Thus the high-distortion
\(a=6\) branch remains inconclusive and is assigned no spectral symmetry
class.  R102 froze a fourth grid, \(h=0.0175\), only for the surviving
\(a=1.02\) core.  The full failure-first sequence is shown in
\cref{fig:quantum-convergence}.

\begin{figure*}[t]
  \centering
  \includegraphics[width=0.98\textwidth]{fig4_quantum_convergence.pdf}
  \caption{Failure-first quantum convergence audit.  The core
  \(a=1.02\) adjacent ratios stabilize across four meshes and
  \(h^2\) extrapolation, whereas the \(a=6\) branches remain sensitive at
  the available resolutions.  The first \(0.04\to0.03\) comparison failed
  the original one-percent level gate; later refinements are shown rather
  than substituted for it.}
  \label{fig:quantum-convergence}
\end{figure*}

\subsection{Converged core window}

On the unified 25--164 window, the R102 \(0.0225\to0.0175\) audit gave
median relative level changes of \(0.347\%\) and \(0.346\%\) for
\(B=0\) and \(1\).  The fine-grid and new \(h^2\)-extrapolated mean ratios
are listed in \cref{tab:core-quantum-window}.  Fine-to-extrapolated ratio
differences are \(0.00242\) at \(B=0\) and \(0.00035\) at \(B=1\); the
correlations between the two pointwise ratio arrays are \(0.960\) and
\(0.974\).  The exact-clock extrapolated mean spacings are \(0.99915\) and
\(0.99903\), providing a finite-window check of the analytic mean density.

\begin{table}[t]
 \centering
 \caption{R102 core results over 140 interior levels.  CDF entries are sup
 distances to the ratio surmises and are descriptive; no \(p\)-values are
 attached.}
 \label{tab:core-quantum-window}
 \small
 \begin{tabular}{@{}lrr@{}}
 \toprule
 Diagnostic & \(B=0\) & \(B=1\) \\
 \midrule
 median level change, \(0.0225\to0.0175\) & 0.347\% & 0.346\% \\
 fine-grid \(\meanratio\) & 0.532251 & 0.587617 \\
 \(h^2\)-extrapolated \(\meanratio\) & 0.529831 & 0.587270 \\
 extrapolated mean unfolded spacing & 0.999151 & 0.999027 \\
 fine CDF distance to GOE & 0.0448 & 0.1140 \\
 fine CDF distance to GUE & 0.1526 & 0.0614 \\
 \bottomrule
 \end{tabular}
\end{table}

The scalar ratio is close to the GOE reference, consistent with the
\(T^2=+1\) symmetry identified in \cref{sec:active}.  At \(B=1\), the mean
and the empirical CDF move toward the GUE reference.  We call these spectra
\emph{GOE-like} and \emph{unitary-leaning} only at this finite-window,
descriptive level.  The 138-ratio CDFs are displayed in
\cref{fig:ratio-cdfs}.

\begin{figure}[t]
  \centering
  \includegraphics[width=0.72\textwidth]{fig5_ratio_cdfs.pdf}
  \caption{Empirical adjacent-ratio distributions over the frozen 138-ratio
  interior window.  The scalar spectrum is closer to the GOE surmise, while
  \(B=1\) moves toward the GUE surmise.  Sup distances are deterministic
  descriptive diagnostics, not i.i.d. goodness-of-fit tests.}
  \label{fig:ratio-cdfs}
\end{figure}

The radial member cannot serve as an untreated Poisson baseline.  On the
R101 fine grid its mean ratio is \(0.045\), and approximately \(39\%\) of
the retained gaps are near-degenerate.  These values reflect the continuum
\(m/-m\) angular-momentum doublets and their finite-grid \(D_4\) remnants.
A symmetry-sector decomposition would be required for a meaningful
integrable-spacing comparison.  Here the radial member is used only as a
degeneracy and classical-integrability control.

\subsection{Frozen magnetic response}

R103 fixed \(B\in\{0,0.25,0.5,1,2,4\}\) before execution and evaluated all
fields on the same \(h=0.0225\) grid and interior window.  The mean-ratio
sequence was
\begin{equation}
 0.5292,\quad0.5439,\quad0.5768,\quad0.5874,\quad0.6229,
 \quad0.5980.
 \label{eq:magnetic-ratio-sequence}
\end{equation}
Every nonzero field lies above the scalar baseline, so the preregistered
retention gate passed.  The response is not monotone.  In particular, the
largest value at \(B=2\) is not promoted as an optimum or assigned arithmetic
meaning.  R104 recomputed the four new fields on the \(h=0.03\) grid; median
level changes were \(0.677\)--\(0.685\%\), and coarse/fine mean-ratio
differences were \(0.0022\)--\(0.0079\), passing both frozen gates.  The
scan therefore supports a resolved field response rather than a single
fortuitous value at \(B=1\), while remaining a low-window crossover rather
than an asymptotic universality result.

\begin{figure*}[t]
  \centering
  \includegraphics[width=0.94\textwidth]{fig6_magnetic_crossover.pdf}
  \caption{Magnetic response on the field grid frozen before execution; the
  uneven field values are placed at categorical horizontal positions for
  readability.  Every nonzero field lies above the scalar baseline, but the
  response is not monotone and no field is selected as an arithmetic optimum.
  Independent coarse/fine pairs for the four new fields retain the direction
  of the response.}
  \label{fig:magnetic-crossover}
\end{figure*}

\subsection{Gauge, solver, and independent-order checks}

R105 repeated the \(a=1.02,h=0.03\) calculation with a deterministic Lanczos
initial vector.  The maximum relative eigen-residual was
\(5.99\times10^{-10}\), and the maximum orthogonality defect was
\(8.04\times10^{-13}\).  The \(B=1\) spectra in symmetric and Landau gauges
agreed to relative \(2.51\times10^{-14}\); the \(B=1\) and \(-1\) spectra
were identical at stored precision; and reruns agreed with the archived R100
arrays within \(3.64\times10^{-14}\).  Matrix-level tests also checked
Hermiticity, Peierls plaquette flux, zero-field reality, and lattice gauge
covariance.

R107 then changed the leading kinetic truncation error by using nearest- and
next-nearest-neighbor Peierls links in a fourth-order covariant finite-
difference stencil.  Its initial run passed every cross-stencil level and
ratio gate but failed the preregistered maximum-residual gate: median Ritz
residuals were already near \(5\times10^{-12}\), while a few edge pairs
reached \(1.9\times10^{-6}\) at \(B=0\) and \(5.0\times10^{-6}\) at
\(B=1\).  The failure remains part of the audit record.

R107A was frozen only after that failure and changed one eigensolver guard:
it requested 200 instead of 180 Ritz pairs, tightened the ARPACK tolerance
from \(2\times10^{-10}\) to \(10^{-12}\), sorted them, and retained the
lowest 180.  No physical parameter, grid, box, field, analysis window,
statistic, or threshold changed.  The retained fine-grid maximum residuals
fell to \(4.65\times10^{-10}\) and \(6.86\times10^{-11}\).  All original
R107 gates then passed, with the quantitative comparison in
\cref{tab:fourth-order-check}.

\begin{table}[t]
 \centering
 \caption{R107A fourth-order finite-difference check.  The reference is the
 archived R102 second-order \(h^2\) extrapolation.}
 \label{tab:fourth-order-check}
 \small
 \begin{tabular}{@{}lrr@{}}
 \toprule
 Diagnostic & \(B=0\) & \(B=1\) \\
 \midrule
 fourth-order coarse/fine median level change & 0.04072\% & 0.04029\% \\
 fourth-order fine vs. second-order extrap. median & 0.01943\% & 0.01894\% \\
 fourth-order fine \(\meanratio\) & 0.534773 & 0.590890 \\
 second-order extrapolated \(\meanratio\) & 0.529831 & 0.587270 \\
 mean-ratio difference & 0.004942 & 0.003620 \\
 pointwise ratio correlation & 0.997432 & 0.995218 \\
 fourth-order mean unfolded spacing & 0.998484 & 0.998364 \\
 max. retained fine residual & \(4.65\times10^{-10}\) & \(6.86\times10^{-11}\) \\
 \bottomrule
 \end{tabular}
\end{table}

R107A is an \emph{independent-order finite-difference check}, not a second
independent discretization family.  Both stencils use a Cartesian grid,
Peierls links, a point-sampled potential, an adapted rectangular wall, and
the ARPACK eigensolver.  Their physical rectangles also differ slightly with
nominal spacing because the geometry includes three mesh steps of padding.
The agreement rules out the simplest explanation based on second-order
kinetic dispersion, but a fixed-domain magnetic finite-element or sine-
Galerkin calculation remains open.

Finally, \cref{fig:smooth-clock} shows the exact-clock unfolding.  Its near-
unit mean spacing checks only the smooth density already proved in
\cref{thm:magnetic-weyl}; the residual staircase fluctuations are not
compared with zeta ordinates.

\begin{figure*}[t]
  \centering
  \includegraphics[width=0.94\textwidth]{fig7_smooth_clock.pdf}
  \caption{Exact-clock unfolding of the extrapolated \(a=1.02\) spectra.
  The cumulative unfolded index follows unit mean spacing for both fields;
  the residual panel displays finite-window fluctuations only.  No Riemann
  ordinate enters this diagnostic.}
  \label{fig:smooth-clock}
\end{figure*}

Across R100--R107A, the defensible conclusion is limited: within 140
interior levels, the centered scalar spectrum is GOE-like and the fixed-field
family exhibits a converged response toward the unitary class.  The data do
not establish GOE/GUE universality, high-energy window stability, or an
endogenous prime-power trace
\citep{BohigasGiannoniSchmit1984,Haake2010,PandeyMehta1983,SaitoEtAl2009}.
No prime or zero arrays were loaded in R100--R107A; the inherited,
zero-exposed parameter lineage is disclosed in \cref{sec:introduction}.
